% listings.tex - Scheme code listing configuration

\lstdefinelanguage{Scheme}{
    % Policy: bold syntax (special forms and library declarations) only.
    % Ordinary procedures (car, map, display, ...) stay unbolded so the
    % highlighting marks what is evaluated specially, not what is common.
    morekeywords={
        define, lambda, if, cond, case, else, and, or,
        when, unless, begin, set!, quote, quasiquote,
        unquote, unquote-splicing,
        let, let*, letrec, letrec*, let-values, let*-values,
        define-values, do, case-lambda,
        define-syntax, syntax-rules, let-syntax, letrec-syntax,
        define-library, import, export, include,
        define-record-type, guard,
        delay, delay-force, parameterize,
        cond-expand, only, except, rename, prefix,
    },
    sensitive=true,
    morecomment=[l]{;},
    morecomment=[s]{\#|}{|\#},
    morestring=[b]{"},
    alsoother={\#},
    literate=
        {->}{$\rightarrow$}1,
}

\lstdefinestyle{scheme}{
    language=Scheme,
    basicstyle=\ttfamily\small,
    keywordstyle=\bfseries\color{keywordcolor},
    commentstyle=\itshape\color{commentcolor},
    stringstyle=\color{stringcolor},
    backgroundcolor=\color{codebg},
    frame=single,
    framerule=0.4pt,
    rulecolor=\color{codeborder},
    xleftmargin=4pt,
    xrightmargin=4pt,
    framexleftmargin=4pt,
    framexrightmargin=4pt,
    aboveskip=\baselineskip,
    belowskip=\baselineskip,
    columns=fullflexible,
    keepspaces=true,
    showstringspaces=false,
    tabsize=2,
    breaklines=true,
    breakatwhitespace=true,
    numbers=none,
    escapeinside={(*@}{@*)},
}

\lstdefinestyle{repl}{
    style=scheme,
    backgroundcolor=\color{outputbg},
    rulecolor=\color{outputborder},
}

\lstset{style=scheme}

% Shorthand for inline code
\newcommand{\code}[1]{\lstinline[style=scheme]{#1}}
